\section{Failure Mode Analysis}
\label{app:failure_modes}

This appendix documents three failure modes encountered during training with
broader implications for RLAIF pipelines: the \emph{Bicameral Mind Problem}
(feedback hallucination in retry generation), the \emph{Retry Assumption
Failure} (inverted preference labels from unconditionally favoring retries),
and the \emph{Group~C Regression} (catastrophic divergence from full-LoRA
training on a Mixture-of-Experts base).

\subsection{The Bicameral Mind Problem}

During iteration~10, we discovered that 35.6\% of the model's thinking
sections contained hallucinated feedback---the model manufactured judge
criticism even when no prior feedback existed in the prompt.

\textbf{Quantification (644 samples):}
\begin{itemize}
    \item 226 contaminated samples from retry attempts
    \item Only 3 contaminated samples from first-attempt generations
    \item 28\% of ``chosen'' preference pairs referenced non-existent feedback
\end{itemize}

\textbf{Example (no feedback in prompt):}
\begin{quote}
\small
``We are now in a correction phase. The judge's feedback is severe: the
agent failed to conduct any searches...'' [No feedback was actually provided]
\end{quote}

\textbf{Root cause:} During training, retry attempts included judge feedback
in the conversation. The model learned to associate the debate task with
receiving criticism, hallucinating this pattern in fresh prompts---a form
of in-context learning contamination.

\textbf{Solution: Proactive self-coaching.} Replace reactive framing
(``the judge said you failed...'') with proactive framing (``when
approaching this task, I should verify every citation...'').
This prevents contamination because proactive thoughts don't require
imagining an external critic, and transfer cleanly to inference.

\textbf{General principle:} Training patterns requiring external
scaffolding will be hallucinated when that scaffolding is absent.

\subsection{The Retry Assumption Failure}

We initially assumed retries were always better than originals, adding
a fixed +0.2 score bonus to all retries. Dual evaluation revealed:

\begin{table}[h]
\centering
\begin{tabular}{lccc}
\toprule
\textbf{Iter} & \textbf{Assumed Gap} & \textbf{Actual Gap} & \textbf{Retries Better} \\
\midrule
1 & +0.183 & $-$0.016 & 48\% \\
2 & +0.193 & $-$0.007 & 49\% \\
3 & +0.191 & $-$0.008 & 48\% \\
4 & +0.167 & $-$0.032 & 44\% \\
\bottomrule
\end{tabular}
\caption{Retry assumption vs. reality. Retries outperform originals only 45\% of the time.}
\end{table}

Cross-examination retries showed dramatic degradation: CX scores dropped
from 0.448 to 0.067--0.233 (60--85\% degradation). Hindsight causes loss
of natural conversational flow.

\textbf{Solution: Dual evaluation.} Score both original and retry
independently. Use higher score as chosen regardless of source. Minimum
gap filter ($\geq 0.10$). Skip retry for scores $\geq 0.85$. Result:
0\% inverted preference pairs (vs. 55\% before).

\subsection{GRPO Regression on MoE Architectures}

Group~C GRPO training on Qwen3-30B-A3B initially produced a catastrophic
regression: mean speech length dropped from $\sim$400 words to $\sim$100
words within three epochs, and validation accuracy fell below random
(45\%). Post-mortem identified three contributing factors. First,
\emph{aggressive LoRA targeting}: we initially adapted all seven linear
projections (q/k/v/o + gate/up/down), producing 3.37B trainable parameters
on a 30B base. For a Mixture-of-Experts model, adapting the expert gating
projections (gate/up/down) destabilizes sparse expert routing under
importance-weighted off-policy updates. Second, \emph{wide epsilon clipping}:
the PPO-style clip parameter $\varepsilon$ was set to $0.3$, permitting
importance ratios to range over $[0.77, 1.30]$ per step. Combined with
a Mixture-of-Experts base, large ratios amplified into large output-space
divergence. Third, \emph{stale logprobs}: because Group~B had merged into
the canonical model between Group~B logprob capture and Group~C training,
the reference policy $\pi_{\text{ref}}$ used in importance ratios was
no longer the policy that generated the samples, producing biased
gradient estimates. The combined effect was runaway divergence.
The fix (Appendix~\ref{app:grpo_groups}) was attention-only LoRA (53M
parameters, a 64$\times$ reduction), tightened $\varepsilon = 0.15$, and
mandatory logprob recomputation after each group merge. These changes
together eliminated the regression with no accuracy loss, and we adopted
them as standard practice for all subsequent groups.
